\structural{ВИСНОВКИ}

У курсовій роботі розроблено імітаційну модель інформаційної системи, яка обробляє завдання
трьох типів, що конкурують за оперативну памʼять, з пʼятьма стратегіями пріоритетного
упорядкування черги. Систему формалізовано стохастичною мережею Петрі з 11~позицій і
7~переходів: обмеження на одночасне виконання завдань відтворено позицією памʼяті й кратністю
дуг, стратегії -- пріоритетами та ймовірностями запуску переходів, правило резервування памʼяті
для завдань вищого пріоритету -- переходом резервування з інформаційною дугою. Алгоритм імітації
реалізовано мовою Java на основі ядра PetriObjModelPaint, класи якого використано без змін, а
потрібні модифікації (відтворюваність прогонів, збір статистики після перехідного періоду,
лічильники подій) внесено в класах-нащадках.

Модель відтворює умову задачі з такими припущеннями: надходження завдань -- пуассонівські
потоки; одночасно виконуються не більше двох завдань A і~B; обслуговування без переривання; пріоритет діє на всю чергу, тож завдання нижчого
пріоритету не обходить завдання вищого пріоритету, що чекає памʼяті; за рівних пріоритетів A і~B
вибір між ними випадковий; памʼять -- єдиний обмежений ресурс, а затримки диспетчеризації
нехтовно малі. Інтенсивності надходження з умови задачі не вдалося дослідити в сталому режимі:
система за них перевантажена ($\rho = 1{,}175$), тому мету моделювання досягнуто в області
стійкості -- за спільного зменшення інтенсивностей коефіцієнтом $\gamma \in [0{,}4;\,0{,}7]$, а
нестійкість за $\gamma = 1$ показано окремо; за $\gamma = 0{,}8$ першими втрачають стійкість
стратегії з найнижчим пріоритетом завдань~C.

Верифікація підтвердила коректність моделі: за $4{,}8\cdot10^8$ подій не порушено ні
інваріант памʼяті, ні правила упорядкування; відгуки змінюються в очікуваному напрямі за зміни
кожної вхідної змінної, а для систем M/M/1 і M/M/2 збігаються з аналітичними (до 2~\%). Експериментальне дослідження (перехідний період
$2\cdot10^4$~хв, час моделювання $3{,}5\cdot10^6$~хв, 60~прогонів у точці, точність 2~\% з
довірчою ймовірністю 0,95, 2100~прогонів факторного плану) показало значущий вплив стратегії на
всі три відгуки за кожної інтенсивності. Стратегії A~$>$~B~$>$~C і B~$>$~A~$>$~C
рівноцінні між собою й найкращі за всіма відгуками: середній час перебування в них у
1,25~раза менший за $\gamma = 0{,}4$ і у 2,0--2,3~раза менший за $\gamma = 0{,}7$, ніж у
стратегіях з пріоритетом~C. Ціна цього -- довге перебування завдань~C (за $\gamma = 0{,}7$ до 101~хв проти 22~хв).
Серед стратегій з пріоритетом~C найкраща C~$>$~A~$>$~B. Залежність відгуків від інтенсивності
описано адекватними за критерієм Фішера рівняннями регресії від нормалізованого фактора
$1/(1-\rho)$ з відносною похибкою не більше 1,2~\%.

Запропоновано збільшити памʼять системи з двох до трьох умовних одиниць -- зміну ресурсу, що
знімає монополізацію завдань~C як наслідок. Пробний запуск
показав, що за потоку з умови задачі система стійка із середнім часом
перебування 15,8--20,5~хв, а за $\gamma = 0{,}7$ час перебування скорочується на 70--86~\%,
черга й очікування -- на 86--94~\%; вибір стратегії стійкості не дає, тож за
потоку, близького до заданого, пропозиція доцільна. Отже, мету моделювання
досягнуто: отримано кількісні рекомендації щодо вибору стратегії та ресурсу.

Імітаційний метод дав змогу дослідити систему, для якої через монополізацію ресурсу й
резервування памʼяті аналітичний розвʼязок отримати складно, і кількісно оцінити пропозицію до
її впровадження. Формалізм стохастичної мережі Петрі виявився зручним для цієї задачі: усі
правила системи задано стандартними елементами мережі (кратністю дуг, пріоритетами,
ймовірностями запуску, інформаційною дугою) без спеціальних алгоритмів, а універсальний алгоритм
імітації мережі не залежить від конкретної системи. Найскладнішим було відтворити правило
резервування, бо пріоритети переходів діють лише на одночасно готові події.
Переваги програми -- відтворюваність кожного прогону за seed і висока швидкість: прогін
тривалістю $3{,}5\cdot10^6$~хв виконується менш ніж за секунду.

Подальше дослідження системи доцільно спрямувати на варіювання інтенсивності окремих типів
завдань у ширшому діапазоні, на гібридні стратегії, у яких пріоритет завдання~C зростає з
довжиною його черги, і на вибір обсягу памʼяті за критерієм вартості.
